\documentclass{beamer}
\usetheme{default}
\usecolortheme{seahorse}

\usepackage{amsmath,amssymb,mathrsfs,bm}
\usepackage{anyfontsize}
\usepackage{unicode-math}
\unimathsetup{math-style=ISO, bold-style=ISO, mathrm=sym}
\setmathfont{XITS Math}
\usepackage{fontspec}
\setmonofont{JetBrains Mono NL}
\AtBeginDocument{
  \let\uglyepsilon\epsilon
  \renewcommand\epsilon{\varepsilon}
}
\AtBeginDocument{
  \let\uglymathsf\mathsf
  \renewcommand\mathsf{\symsf}
}
\renewcommand{\bm}{\symbf}

\usepackage[fontset=none]{ctex}
\setCJKmainfont{Adobe Song Std}[BoldFont=Noto Sans CJK SC Bold,ItalicFont=Adobe Kaiti Std]
\setCJKsansfont{Adobe Song Std}[BoldFont=Noto Sans CJK SC Bold,ItalicFont=Adobe Kaiti Std]
\setCJKmonofont{Noto Sans Mono CJK SC}[RawFeature=+fwid]
\newCJKfontfamily\songti{Adobe Song Std}[BoldFont=Noto Sans CJK SC Bold]
\newCJKfontfamily\heiti{Noto Sans CJK SC}[RawFeature=+fwid]
\newCJKfontfamily\kaishu{Adobe Kaiti Std}
\newCJKfontfamily\fangsong{Adobe Fangsong Std}

\setbeamerfont{title}{series=\bfseries}
\setbeamerfont{author}{family=\kaishu}
\setbeamerfont{date}{family=\kaishu}
\setbeamertemplate{navigation symbols}{}
\setbeamertemplate{items}[ball]
\setbeamertemplate{section in toc}[ball unnumbered]
\setbeamertemplate{footline}[frame number]
\setbeamerfont{section in toc}{series=\bfseries}
\setbeamertemplate{part page}{
        \begin{beamercolorbox}[sep=8pt,center,wd=\textwidth]{part title}
            \bfseries\Large\insertpart\par
        \end{beamercolorbox}
}
\renewcommand\baselinestretch{1.25}

\usepackage{multicol}
\setlength{\columnsep}{25pt}
\usepackage{color}
\usepackage{graphicx,caption,subcaption}
\usepackage{hyperref}
\hypersetup{
    colorlinks=true,
    linkcolor=blue,
    filecolor=green,
    urlcolor=cyan,
    citecolor=cyan,
}

\usepackage{listings}
\lstset{
    language=Python,
    basicstyle=\ttfamily\tiny,
    aboveskip={1.0\baselineskip},
    belowskip={1.0\baselineskip},
    xleftmargin=25pt,
    columns=fixed,
    extendedchars=true,
    breaklines=true,
    tabsize=4,
    prebreak=\raisebox{0ex}[0ex][0ex]{\ensuremath{\hookleftarrow}},
    frame=lines,
    showtabs=false,
    showspaces=false,
    showstringspaces=false,
    keywordstyle=\color[rgb]{0,0.5,0},
    commentstyle=\color[rgb]{0.5,0,0},
    stringstyle=\color[rgb]{1,0,0},
    numbers=left,
    numberstyle=\tiny\ttfamily,
    stepnumber=1,
    numbersep=12pt,
    captionpos=t,
    escapeinside={\%*}{*)}
}

\title{我的Beamer模板}
\author{作者}
\date{\today}

\begin{document}

\begin{frame}
    \titlepage
\end{frame}

\part{第一部分}

\begin{frame}
    \partpage
    \tableofcontents
\end{frame}

\section{第一节}

\begin{frame}{Frame Title 1}
    \begin{center}
        赴戍登程口占示家人二首（其二）
        \medskip

        \fangsong{【清】林则徐}
        \bigskip

        \kaishu{
        力微任重久神疲，再竭衰庸定不支。

        苟利国家生死以，岂因祸福避趋之！

        谪居正是君恩厚，养拙刚于戍卒宜。

        戏与山妻谈故事，试吟断送老头皮。
        }
    \end{center}
\end{frame}

\section{第二节}

\begin{frame}{Frame Title 2}
    \begin{align*}
        [q,p]_{Q,P}
        &= \frac{\partial q}{\partial Q}\cdot\frac{\partial p}{\partial P} - \frac{\partial q}{\partial P}\cdot\frac{\partial p}{\partial Q} \\
        &= \left(-\frac{\sqrt{1-P^2\symrm{e}^{2Q}}}{\symrm{e}^Q} - \frac{P^2\symrm{e}^Q}{\sqrt{1-P^2\symrm{e}^{2Q}}}\right)\cdot\left(\frac{-\symrm{e}^Q}{\sqrt{1-P^2\symrm{e}^{2Q}}}\right) - \\
        &\quad\left(\frac{-P\symrm{e}^Q}{\sqrt{1-P^2\symrm{e}^{2Q}}}\right)\cdot\left(\frac{-P\symrm{e}^Q}{\sqrt{1-P^2\symrm{e}^{2Q}}}\right) \\
        &= 1 + \frac{P^2\symrm{e}^{2Q}}{1-P^2\symrm{e}^{2Q}} - \frac{P^2\symrm{e}^{2Q}}{1-P^2\symrm{e}^{2Q}} \\
        &= 1
    \end{align*}
\end{frame}

\begin{frame}{Frame Title 3}
    \begin{equation}
    D_n = \frac{2}{l}\int_0^l\frac{2blx^3-abx^4}{12a^2}\sin\frac{n\pi}{2l}x \,\mathrm{d}x = \left\{
    \begin{aligned}
        0\qquad &,\quad n = 2k \\
        -\frac{8bl^4}{a^2 \pi^5 n^5}&,\quad n = 2k + 1
    \end{aligned}
    \right.
    \end{equation}
    \begin{equation}
    w(x,0) = 0,\qquad \left.\frac{\partial w}{\partial t}\right|_{t=0} = -\left.\frac{\partial v}{\partial t}\right|_{t=0} = -\frac{Aa\sin\left(\dfrac{\omega x}{a}\right)}{ES\cos\left(\dfrac{\omega l}{a}\right)}
    \end{equation}
    \begin{equation}
    \iiint \nabla\cdot\symbf{A}\,\symrm{d}V = \oiint \symbf{A}\,\symrm{d}\symbf{S}
    \end{equation}
\end{frame}

\part{第二部分}

\begin{frame}
    \partpage
    \tableofcontents
\end{frame}

\section{第三节}

\begin{frame}{Frame Title 4}
    \begin{itemize}
        \item 我能吞下玻璃而不伤身体
        \begin{itemize}
            \item[1.] \songti{我能吞下玻璃而不伤身体。，、——：；‘“恭喜发财”’！？}

            \item[2.] \kaishu{我能吞下玻璃而不伤身体。，、——：；‘“恭喜发财”’！？}

            \item[3.] \fangsong{我能吞下玻璃而不伤身体。，、——：；‘“恭喜发财”’！？}

            \item[4.] \heiti{我能吞下玻璃而不伤身体。，、\symbol{"2E3A}：；‘“恭喜发财”’！？}
        \end{itemize}

        \item \href{https://www.php.net/}{PHP}是世界上最好的语言

        \item \songti{欢迎来到\textbf{\textcolor[RGB]{139,0,18}{北京大学}}物理学院！}
    \end{itemize}
\end{frame}

\begin{frame}{Frame Title 5}
    \begin{multicols}{2}
    Lorem ipsum dolor sit amet, consectetur adipiscing elit, sed do eiusmod tempor incididunt ut labore et dolore magna aliqua. Ut enim ad minim veniam, quis nostrud exercitation ullamco laboris nisi ut aliquip ex ea commodo consequat. Duis aute irure dolor in reprehenderit in voluptate velit esse cillum dolore eu fugiat nulla pariatur. Excepteur sint occaecat cupidatat non proident, sunt in culpa qui officia deserunt mollit anim id est laborum.
    \end{multicols}
\end{frame}

\section{第四节}

\begin{frame}[fragile]{Frame Title 6}
    模板代码：
    \begin{lstlisting}
    def permutations(a, arr, pos, end):  # 获取列表的全排列
        if pos == end:  # 递归结束
            arr = arr[:]
            a.append(arr)
        else:
            for index in range(pos, end):
                arr[index], arr[pos] = arr[pos], arr[index]  # 交换相邻两个元素
                permutations(a, arr, pos+1, end)  # 列表长度-1
                arr[index], arr[pos] = arr[pos], arr[index]  # 交换相邻两个元素
        return a
    \end{lstlisting}
\end{frame}

\begin{frame}{Frame Title 7}
    Itemize example: \pause
    \begin{itemize}
        \item Ceasar said: “Veni, Vidi, Vici.” \pause
        \item Karl Marx said: “The history of all hitherto existing society is the history of class struggles.” \pause
    \end{itemize}
    Enumerate example: \pause
    \begin{enumerate}
        \item \LaTeX{} beamer is powerful. \pause
        \item \LaTeX{} beamer is beautiful. \pause
    \end{enumerate}
\end{frame}

\begin{frame}{Frame Title 8}
    \begin{columns}
        \begin{column}{0.5\linewidth}
            \colorbox{red}{Fourier变换}
            \[
            F(\omega) = \frac{1}{\sqrt{2\pi}}\int_{-\infty}^{\infty}f(t)\mathrm{e}^{-i\omega t}\,\mathrm{d}t
            \]
        \end{column}
        \begin{column}{0.5\linewidth}
            \colorbox{yellow}{Fourier逆变换}
            \[
            f(t) = \frac{1}{\sqrt{2\pi}}\int_{-\infty}^{\infty}F(\omega)\mathrm{e}^{i\omega t}\,\mathrm{d}\omega
            \]
        \end{column}
    \end{columns}
\end{frame}

\begin{frame}
    \begin{center}
        \Huge\textbf{谢谢大家！}
    \end{center}
\end{frame}

\end{document}
